% !TeX root = ../colloquium-2.tex

% Определение производной. 
% Связь дифференцируемости и производной.
% Непрерывность дифференцируемой функции.

\phantomsection\label{sec_Definition06}
\begin{definition}[Определение производной.
                   Связь дифференцируемости и производной.
                   Непрерывность дифференцируемой функции.]
    \subsection*{Определение производной.}
    \begin{note}[purple]
        Если существует предел 
        $\displaystyle{\lim_{h \rightarrow 0}} \frac{f(a + h) - f(a)}{h}$, 
        то он называется производной функции $f$ в точке $a$. \\
        Другая форма записи предела:
        $$ \lim_{x \rightarrow a} \frac{f(x) - f(a)}{x - a} $$
        Производная функции $f$ в точке $a$ обозначается символом 
        $f'(a)$ или $\frac{df}{dx}(a)$. 
    \end{note}

    \subsection*{Связь дифференцируемости и производной.}
    \begin{note}[purple]
        Функция $f$ дифференцируема в точке $a$ тогда и только тогда, когда 
        в точке $a$ существует производная этой функции $f'(a)$. При этом 
        $df(h) | _{x = a} = f'(a)h$.
    \end{note}

    \subsection*{Непрерывность дифференцируемой функции.}
    \begin{note}[purple]
        Пусть функция $f$ дифференцируема в точке $a$. Тогда $f$ непрерывна в точке $a$. 
    \end{note}
\end{definition}